\originalchapter{8.06 作业3：精细结构与向量算符}
来源为Spring 2018 Assignment 3, 2018年2月21日发布. 本轮第3--5题有完整MIT内容, 中文转述及\textbf{AI独立解答}. 第1题明确打印Griffiths 6.15、p.270“minor modifications”题面, 第2题署Sakurai具体题, 单核真正复制对象; 不将两题恢复为独立MIT题. 原件已归档, 此局部状态不称确证侵权.
\originalsection{第3题：3P精细结构的磁场尺度}
\begin{exercise}
(i)求使自由电子Zeeman间距等于氢$3P_{1/2},3P_{3/2}$分裂的磁场. (ii)估计3P电子内部自旋轨道有效磁场.
\end{exercise}
\begin{solution}
精细能量公式$\delta E_{nj}=-mc^2\alpha^4[1/(j+1/2)-3/(4n)]/(2n^3)$, 两$j$间距$\Delta E=mc^2\alpha^4/(4n^3)$, $n=3$约$1.342\times10^{-5}$eV. 自由电子$g\simeq2$的两自旋间距$2\mu_BB$, 所以$B=\Delta E/(2\mu_B)\simeq0.1159$T.

内部场来自电子轨道运动下电场的相对论变换, 自旋轨道Hamilton算符须含Thomas因子1/2:$H_{\rm so}=e^2L\cdot S/(2m^2c^2r^3)$（Gaussian单位）. 对$n=3,\ell=1$, $\langle r^{-3}\rangle=1/[a_0^3n^3\ell(\ell+1/2)(\ell+1)]=1/(81a_0^3)$. 若以轨道大小$|L|=\sqrt2\hbar$估计, 与电子磁耦合$2\mu_BS\cdot B_{\rm eff}/\hbar$对照给$B_{\rm eff}=e^2\hbar|L|\langle r^{-3}\rangle/(4\mu_Bm^2c^2)\simeq0.1093$T. 若“内部场”指Thomas修正前运动参考系场则约0.2186T. 场是沿$L$的算符相关有效量, 不能把它误当给定方向的经典常场, 两估计与(i)的能级间距尺度同阶而不要求相等.
\end{solution}
\originalsection{第4题：CSCO与强场Zeeman}
\begin{exercise}
(a)限氢束缚态含电子自旋, 判断是否CSCO: i.$H_0,L^2,L\cdot S,J_z$; ii.$H_0,L^2,L_z,S_x$; iii.$H_0,L^2,J_z,S_z$; iv.$H_0,J^2,J_z,S_z$; v.$H_0,J^2,L\cdot S,J_z$. (b)强場未耦合基求$\langle1/(\hat j+1/2)\rangle$, 用$\langle J^2\rangle$确定两个$j$概率, 再求精细位移; 为何尚有$m_\ell+2m_s$简并仍可这样算?
\end{exercise}
\begin{solution}
(i)由$J^2=L^2+3\hbar^2/4+2L\cdot S$, 等价标准$(n,\ell,j,m)$CSCO. (ii)四算符相容且确定$(n,\ell,m_\ell,m_x)$, 是CSCO, 虽不是$S_z$基. (iii)$J_z-S_z=L_z$, 等价未耦合CSCO. (iv)$[J^2,S_z]=2[L\cdot S,S_z]$一般非零, 不是相容组, 特例最高态不改变全局结论. (v)由$L^2=J^2-3\hbar^2/4-2L\cdot S$等价标准组, 是CSCO.

对固定$\ell\ge1$, $\langle J^2\rangle/\hbar^2=\ell(\ell+1)+3/4+2m_\ell m_s$. 两读数$j_+=\ell+1/2,j_-=\ell-1/2$给$p_+=(\ell+1+2m_\ell m_s)/(2\ell+1)$,$p_-=(\ell-2m_\ell m_s)/(2\ell+1)$. 所以
\[
\left\langle\frac1{\hat j+1/2}\right\rangle
=\frac{p_+}{\ell+1}+\frac{p_-}{\ell}
=\frac{\ell(\ell+1)-m_\ell m_s}{\ell(\ell+1/2)(\ell+1)}.
\]
精细位移将此代入$-mc^2\alpha^4[\langle1/(\hat j+1/2)\rangle-3/(4n)]/(2n^3)$, 全Zeeman另加$\mu_BB(m_\ell+2m_s)$. $\ell=0$只有$j=1/2$, 不使用分母为0的中间式, 期望为1.

同一$n$有相同$m_\ell+2m_s$的简并仍需查扰动内部矩阵. 精细Hamilton算符与$L^2,J_z$对易; 不同$\ell$或不同$m_j=m_\ell+m_s$没有非对角元. 同一$\ell$下若Zeeman标签相同且$m_s$不同, $m_\ell$相差2, $m_j$相差1, 也不耦合; 若$m_s$相同则$m_\ell$同而为同一态. 因而该简并块在这些基矢间已对角, 可用对角期望, 不能仅说“强场所以非简并”.
\end{solution}
\originalsection{第5题：向量算符恒等式}
\begin{exercise}
若$[J_i,W_j]=\ii\hbar\epsilon_{ijk}W_k$, $W$称向量算符. (a)证位置对轨道$L$是向量. (b)证两向量叉积仍向量. (c)证$[J^2,W]=2\ii\hbar W\times J+2\hbar^2W$, 核$W=J$. (d)求常数$\alpha$使$[J^2,[J^2,V]]/\alpha=(V\cdot J)J-\{J^2,V\}/2$.
\end{exercise}
\begin{solution}
(a)$[\epsilon_{iab}x_ap_b,x_j]=\epsilon_{iab}x_a(-\ii\hbar\delta_{bj})=\ii\hbar\epsilon_{ijk}x_k$. (b)用乘积对易规则, 两项$\epsilon$缩并与作业10练习6相同, 保持向量顺序, 得$[J_i,(U\times V)_j]=\ii\hbar\epsilon_{ijk}(U\times V)_k$.

(c)$[J^2,W_j]=\sum_i(J_i[J_i,W_j]+[J_i,W_j]J_i)$. 将第一项$J_iW_k$换成$W_kJ_i+[J_i,W_k]$, 两个同序项给$2\ii\hbar(W\times J)_j$, 额外$\epsilon$缩并给$2\hbar^2W_j$. $W=J$时$J\times J=\ii\hbar J$, 两项正好抵消为0.

(d)记$C=[J^2,V]$, 它也向量. 再用(c)给$[J^2,C]=2\ii\hbar C\times J+2\hbar^2C$. 代$C=2\ii\hbar V\times J+2\hbar^2V$, 得$-4\hbar^2(V\times J)\times J+8\ii\hbar^3V\times J+4\hbar^4V$. 有序三重积$(V\times J)\times J=J(V\cdot J)-VJ^2+\ii\hbar V\times J$, 而向量标量$V\cdot J$与每个$J_i$对易; 将$V\times J$用(c)换回并归并可得
\[
[J^2,[J^2,V]]=-4\hbar^2\left[(V\cdot J)J-\tfrac12(J^2V+VJ^2)\right].
\]
因此$\alpha=-4\hbar^2$. 该式可在两个角动量子空间间取矩阵元独立核对, $W=J$亦使两侧0, 避免将经典叉积次序当算符次序.
\end{solution}
